\documentclass[12pt]{article}
\usepackage{geometry}
\usepackage{xcolor}
\usepackage{hyperref}
\usepackage[utf8]{inputenc}
\usepackage{tabto}
\usepackage{graphicx}
\usepackage{biblatex}
\addbibresource{bibliography.bib}
\graphicspath{ {./images/} }
\hypersetup{
colorlinks=false,
linktoc=all,
linkcolor=white,
}
\geometry{
a4paper,
left=30mm,
top=30mm,
right=30mm,
bottom=40mm
}
\title{Animation Vectorization and Compression}
\author{Bassam Helal}
\date{ September 2020 }
\begin{document}
    % TODO 12-Sep-20 Remove black background and white text
    \pagecolor{black}
    \color{white}

    \pagenumbering{gobble}

    \maketitle

    \begin{center}
        \vspace{8cm}
        \includegraphics[scale=0.65]{SwanseaUniversityLogo.png}
    \end{center}

    \pagebreak

    \begin{center}
        \section*{Abstract}
    \end{center}

    % TODO 12-Sep-20 Abstract

    \pagebreak

    \renewcommand*\contentsname{
    \begin{center}
        Table of Contents
    \end{center}}

    \tableofcontents

    \pagebreak

    \pagenumbering{arabic}


    \section{Introduction}

    A paragraph or so for each of:

    \begin{itemize}
        \item Project Description
        \item Motivation
        \item Existing Literature
        \item Limitations of the literature that we will attempt to tackle
        \item Aims and Objectives
        \item Discoveries and results we found after our attempts
        \item Section Signposting
    \end{itemize}

    \pagebreak


    \section{Background}

    \tab
    This document makes no assumption of the reader's knowledge in computer graphics and its related fields,
    thus, this section will introduce the reader to some of the required knowledge needed for understanding this
    project.
    Firstly, we will introduce the concept of Raster Graphics, ie, the current common standard for graphics data
    representation.
    We will then introduce its counterpart Vector Graphics, a different way to store graphics data with its own
    advantages and disadvantages that is growing in popularity.
    We will then give a brief primer on Image compression and the different methods of reducing the overall size
    of an image.
    Finally, we will explain how this all ties into video graphics and describe some of the few differences between
    concepts from image graphics and video graphics.

    \subsection{Raster Graphics}

    1 - 2 sentences for each of:

    \begin{itemize}
        \item Brief summary of what we will go over (signposting)
        \item Pixels, Bitmaps, Image Data Buffers: basically the bare must knows about Raster Graphics
        \item Advantages and popularity of Raster Graphics
        \item Disadvantages like upscaling, pixelation, size issues with high resolution
        \item Brief summary of everything we just mentioned
    \end{itemize}

    \subsection{Vector Graphics}

    1 - 2 sentences for each of:

    \begin{itemize}
        \item Brief summary of what we will go over (signposting)
        \item Declarative vs Imperative ie WYSIWYG vs WYSIWYM, describing the image using primitive geometry
        \item Advantages of Vector especially in contrast to Raster (size and scalability), mention uses and popularity
        \item Disadvantages like portability and light technical stuff like gradients and photographs
        \item Brief summary of everything we just mentioned
    \end{itemize}

    We can make a small table comparing Raster \& Vector Graphics for the attributes like size, scalability etc

    \subsection{Image Compression}

    1 - 2 sentences for each of:

    \begin{itemize}
        \item Brief summary of what we will go over (signposting)
        \item Fundamentals of Compression (type-agnostic) things like compression rate and lossy-ness or accuracy
        \item Lossless Image Compression (PNG, BMP etc)
        \item Lossy Image Compression (JPEG, GIF etc)
        \item Brief summary of everything we just mentioned
    \end{itemize}

    We can have an image comparing popular image file types with regards to compression

    \subsection{Video Graphics}

    This section may be quite short

    \begin{itemize}
        \item Brief summary of what we will go over (signposting)
        \item How video is stored and how it relates to Images (mostly the same)
        \item Common file types and how they compress and store video data (MP4, MKV, H265 etc)
        \item Evidence and proof of why video is expensive (and also popular)
        \item Brief summary of everything we just mentioned
    \end{itemize}

    \pagebreak


    \section{Related Work}

    Here we throw a lot of the current stuff from both research and industry and show their flaws and limitations
    (which we may or may not have tackled)

    \subsection{Current Research}

    \begin{itemize}
        \item Brief summary of what we will go over (signposting)
        \item New Machine Learning methods for Vectorization
        \item New Machine Learning methods for Compression
        \item Novel approaches and ideas within Vectorization such as mesh gradients etc (no ML)
        \item Brief summary of everything we just mentioned
    \end{itemize}

    \subsection{Current Applications}

    \begin{itemize}
        \item Brief summary of what we will go over (signposting)
        \item Adobe Illustrator, Logos, Topography, Medical stuff like X rays etc Potrace
        \item Animations and Games using Vector as the development method
        \item Current Image \& Video compression algorithms such as PNG, H265 and H266 (no ML)
        \item Brief summary of everything we just mentioned
    \end{itemize}

    \subsection{Limitations}

    \begin{itemize}
        \item Brief summary of what we will go over (signposting)
        \item Why ML is absolute hot trash
        \item Show off the flaws with existing technologies while showing how we tackled (and hopefully beat) them
        \item Discuss inherent limitations of this process, such as issues with Vector Graphics and input data being varying etc
        \item Brief summary of everything we just mentioned
    \end{itemize}

    \pagebreak


    \section{Implementation}

    This may get long and technical, we should try to keep things fairly high level as much as possible,
    we will also show both Math, pseudocode and pipeline of the overall process and avoid boring low level details

    \subsection{Pipeline}

    The very basic common high level pipeline demonstrating the input, the processing, and the output

    \subsubsection{Input}

    An animated cartoon video, particularly one with vectorizable qualities, which we will further detail
    \begin{itemize}
        \item Low or no use of gradients so it has discrete colors
        \item Low or no noise
        \item No 3D stuff like shadows and reflections
    \end{itemize}

    \subsubsection{Processing}

    We have tried multiple processes which we further explain in the approaches section, but here
    we should go over the general idea of what is necessary so:
    \begin{itemize}
        \item Frame extraction
        \item Conversion of each Frame to an SVG which will have a single SVG Path element for each color
        \item SVG Frame Joining which will have some smart compression built-in
    \end{itemize}

    \subsubsection{Output}

    An SVG Video that is of course highly scalable and has a lower size than the original if upscaled and with
    very minimal data loss

    \subsection{Toolchain}

    No need to spend too much time on this part, just quickly mention the tools that we use and why we chose them.
    This makes the reader brace themselves for what we did without actually telling them just yet.

    \subsection{Approaches}

    Here we will explain our story of the last few months, the many approaches and their failures. This section
    should implicitly show the reader that the problem is very hard, harder than initially thought to be. If
    we show the difficulty and the struggle it can explain any poor results and can further show overall contribution
    that we tried something difficult

    \subsubsection{Color Quantization Approach}

    From the general process of the library ImageTracer.js

    \begin{itemize}
        \item Explain Color Quantization
        \item Advantages (very accurate, fast, good for logos and web based stuff)
        \item Disadvantages (Too dependent on quantization number of colors, large size and poor speed with high number of colors)
        \item Show pictures of original vs some different results from this approach (to show why it fails)
    \end{itemize}

    \subsubsection{Connected Component Labelling (CCL)}

    Using Connected Component Labelling (CCL) to create discrete but adjacent components

    \begin{itemize}
        \item Explain CCL
        \item Reasoning behind approach
        \item Possible disadvantages
        \item Signpost for the following approaches
    \end{itemize}

    \subsubsection{Edge Based CCL Approach}

    Using OpenCV Canny Edge detection

    \begin{itemize}
        \item Explain OpenCV and Canny Edge Detection
        \item Reasoning behind approach (edges should form polygons which will be color regions) and the expected result
        \item Advantages (fast, easy, accurate for edge detection)
        \item Disadvantages (not all edges form polygons and 1 pixel incorrect can have huge effects)
    \end{itemize}

    \subsubsection{Pixel Based CCL Approach}

    Using the color of each pixel to form regions with close enough or similar colors

    \begin{itemize}
        \item Explain pixel data (RGBA) and what it means to be close to a neighboring pixel
        \item Reasoning behind approach (an image can be a partition of several color regions)
        \item Advantages (surprisingly accurate, little data loss, huge size savings)
        \item Disadvantages (slow, complex, Anti-Aliasing and noise affects result so too many 1 pixel regions)
    \end{itemize}

    \subsection{Limitations \& Drawbacks}

    Here we explain the overall limitations of all the processes we used and really any future process because
    of the nature of Vectorization of Raster data

    \begin{itemize}
        \item Brief summary of what we will go over (signposting)
        \item Speed or performance (never realtime and requires HPC on GPUs for good speeds) and optimizations required
        \item Lossy-ness, the transformation of Raster to Vector is by nature lossy because we are undoing Rasterization
        artifacts
        \item Size, this can be an issue for low resolution images or for very complex images, rasterization is often better
        \item Brief summary of everything we just mentioned
    \end{itemize}

    \pagebreak


    \section{Software Engineering}

    Just the SE stuff and how we applied it during development, briefly mention the changes Covid did to the development

    \subsection{Methodology}

    Extreme Lean Agile because life is unknown beyond 2 days in these strange times we live in :/
    \linebreak
    Here we can show how the poor results affected our schedule and how things shifted because of unexpected results,
    we can definitely show that this is common in real life and is a good representation and preparation of the real world
    which has unexpected delays and poor results all the time and how we overcame it all (less sleep and more coffee :D)

    \subsection{Schedule (Gantt Chart)}

    Followed it well but started falling apart towards the end, very accurate to the real world :D
    \linebreak
    Mention that we started dissertation somewhat early to avoid the risk of running out of time

    \subsection{Risks}

    Time Time Time!!!

    \subsection{Testing?}

    Do we need this section?
    \linebreak
    Unit Testing for accuracy and of course Manual Acceptance testing (to ensure images actually open, are not corrupt etc)

    \pagebreak


    \section{Evaluation}

    Two main angles, accuracy to the original and compression rate from the original
    \linebreak
    As much as possible we need to emphasize that we have done something great,
    no lying, just selective and sometimes exaggerated benchmarks in order to further our claim that we
    have accomplished something meaningful

    \subsection{Accuracy}

    \begin{itemize}
        \item Define accuracy (aka lossy-ness or how much data is lost from original)
        \item Quantitative methods (compare each pixel to itself), and compare with other compression methods like JPEG
        \item Qualitative methods (human eye side by side) and compare with other methods like JPEG
    \end{itemize}

    \subsection{Compression}

    \begin{itemize}
        \item Define compression rate (aka how smaller is the result compared to original or some other)
        \item Quantitative methods (compare against original, AND against others including upscaled and uncompressed)
        \item Qualitative methods (file transfer speeds of smaller sizes and how smaller files are better for users)
    \end{itemize}

    \pagebreak


    \section{Results \& Findings}

    Here we show both the results of our development but also our own personal findings that someone who
    wants to research this kind of thing should be aware of that I discovered

    \subsection{Results}

    TODO
    \linebreak
    The good, the bad and the ugly (with detailed explanations :D)

    \subsection{Findings}

    \begin{itemize}
        \item This is a very difficult problem, more difficult than initially thought to be
        \item GPU acceleration is very hard but improves speed shockingly well
        \item Node sucks for concurrency and doesn't have true multi-threading
        \item Differing input will have different results because of the limitations of vectorization
    \end{itemize}

    \pagebreak


    \section{Challenges}

    \begin{itemize}
        \item Toolchain used was not very mature or helpful sometimes
        \item GPU programming is insanely difficult but yields huge benefits (leading to very difficult decisions)
        \item Optimizations to increase speed are also difficult and especially so on a non-multi-threaded environment like Node
        \item Reducing the number of human inputted arguments to create a fully autonomous deterministic system that is not ML based and produces very accurate results is very difficult
    \end{itemize}

    \pagebreak


    \section{Reflection}

    \begin{itemize}
        \item How would you do it differently if to start from scratch (Use JVM or Native and optimize for GPU early)
        \item Project Management (Fairly good given the circumstances)
        \item So much new knowledge about many different things
        \item Very difficult but very mentally stimulating and rewarding project that I really enjoyed
    \end{itemize}

    \pagebreak


    \section{Future Work}

    \begin{itemize}
        \item Use better tools that allow for native multi-threading and easier and direct access to GPU
        \item Better optimizations and better compression
        \item Explore having a mixed raster and vector solution to solve the areas where vector fails
        \item Buy better hardware for faster development :D
    \end{itemize}

    \pagebreak


    \section{Conclusion}

    \begin{itemize}
        \item Summarize everything we said
        \item Conclude with final thoughts about the good, the bad and how this can further improve the world etc
    \end{itemize}

    \pagebreak


    \printbibliography[heading=bibintoc,title={References}]

    \pagebreak

\end{document}